\chapter{Đặc tả yêu cầu}\label{ch:yeucau}
Đặc tả đầy đủ nằm ở \path{docs/requirements/srs.md}, \path{context-and-scope.md} và \path{use-case-checkout.md}. Hệ mã FR/NFR được giữ xuyên suốt để lập ma trận truy vết ở \cref{ch:truyvet}.

\section{Bối cảnh, phạm vi (trong / ngoài)}\label{sec:se-phamvi}
\textbf{Vấn đề nghiệp vụ.} Quán cà phê quy mô nhỏ cần để khách đặt món và thanh toán không dùng tiền mặt trên điện thoại, tránh xếp hàng và nhầm đơn.

\textbf{Trong phạm vi MVP:} xem menu và chi tiết (size S/M/L, topping); giỏ hàng; đăng ký, đăng nhập, JWT, phân quyền CUSTOMER và STAFF; tạo đơn; thanh toán Ví hoặc Thẻ mô phỏng; xác nhận, theo dõi trạng thái, lịch sử đơn; máy trạng thái, thanh toán idempotent, kiểm soát tồn kho đồng thời, khuyến mãi, điểm thưởng; triển khai production qua CI/CD.

\textbf{Mở rộng (nếu còn thời gian):} PWA đầy đủ (offline, thông báo); giao diện quản trị; giao diện riêng cho nhân viên quầy. \textbf{Không thực hiện:} cổng thanh toán thật, ứng dụng native, giao hàng và nhiều chi nhánh, quản lý nguyên liệu chi tiết.

\section{Stakeholder}\label{sec:se-stakeholder}
\begin{longtable}{@{}p{3cm}p{3cm}p{4.6cm}p{3.6cm}@{}}
\caption{Stakeholder}\label{tab:stakeholder}\\
\toprule \textbf{Stakeholder} & \textbf{Vai trò} & \textbf{Mục tiêu và kỳ vọng} & \textbf{Ảnh hưởng đến yêu cầu} \\ \midrule \endfirsthead
\toprule \textbf{Stakeholder} & \textbf{Vai trò} & \textbf{Mục tiêu và kỳ vọng} & \textbf{Ảnh hưởng đến yêu cầu} \\ \midrule \endhead
\bottomrule \endlastfoot
Khách hàng & Đặt và thanh toán đơn & Đặt nhanh, giá rõ, theo dõi được đơn & Menu, giỏ, thanh toán, lịch sử \\
Nhân viên quầy & Pha chế, cập nhật trạng thái & Nhận đúng đơn đã thanh toán & \texttt{PATCH /orders/:id/status} \\
Quản lý cửa hàng & Quản lý sản phẩm, khuyến mãi & Kiểm soát tồn kho, khuyến mãi & Tồn kho, khuyến mãi; quản trị ngoài MVP \\
Giảng viên & Khách hàng nghiệm thu & Chạy được, có tài liệu và bằng chứng & Toàn bộ sản phẩm bàn giao \\
Nhóm phát triển & Xây dựng sản phẩm & Đúng hạn, dễ bảo trì & Ràng buộc công nghệ, CI/CD \\
Đơn vị thanh toán (mô phỏng) & Xử lý Ví/Thẻ & Nhận yêu cầu hợp lệ, không trừ tiền lặp & Idempotency, trạng thái thanh toán \\
\end{longtable}

\section{Giả định, ràng buộc, phụ thuộc}\label{sec:se-rangbuoc}
\textbf{Giả định.} (G1) Thanh toán mô phỏng, có thể thành công hoặc thất bại theo kịch bản. (G2) Một cửa hàng, một loại tiền (VND). (G3) Phụ phí size: S +0, M +5.000, L +10.000 đồng. (G4) Điểm thưởng: 1 điểm cho mỗi 10.000 đồng. (G5) Tồn kho chỉ tính trên \texttt{Product.stock}, trừ khi tạo đơn, hoàn lại khi hủy. (G6) Đơn mới ở trạng thái PENDING.

\textbf{Ràng buộc.} (R1) Bắt buộc Next.js, NestJS, PostgreSQL. (R2) Frontend mobile-first/PWA. (R3) Hạn nộp 27/10/2026. (R4) Nhóm năm người, làm theo Scrum. (R5) Production qua HTTPS, CSDL không public, secret không nằm trong repo.

\textbf{Rủi ro ban đầu.} Không kịp triển khai production (giảm thiểu: triển khai thử sớm, 23/10); hai bài tập lớn chồng lịch (dùng chung mã nguồn); đặt đồng thời gây âm tồn kho (cập nhật có điều kiện trong transaction, có test); sơ đồ mâu thuẫn SRS (rà bằng ma trận truy vết); thành viên chậm tiến độ (quy trình trong thỏa thuận nhóm).

\section{Yêu cầu chức năng}\label{sec:se-fr}
Mức ưu tiên theo MoSCoW. FR-16 (Could) và FR-17 (Won't) nằm ngoài MVP.
{\small
\begin{longtable}{@{}p{1.2cm}p{5.4cm}p{1.2cm}p{6.6cm}@{}}
\caption{Yêu cầu chức năng và tiêu chí chấp nhận}\label{tab:fr}\\
\toprule \textbf{Mã} & \textbf{Yêu cầu} & \textbf{Ưu tiên} & \textbf{Tiêu chí chấp nhận} \\ \midrule \endfirsthead
\toprule \textbf{Mã} & \textbf{Yêu cầu} & \textbf{Ưu tiên} & \textbf{Tiêu chí chấp nhận} \\ \midrule \endhead
\bottomrule \endlastfoot
FR-01 & Hệ thống hiển thị danh sách sản phẩm có tên, giá, ảnh (nếu có) & Must & \texttt{GET /\allowbreak{}products} trả danh sách; màn hình menu hiển thị đủ sản phẩm có trong seed \\
FR-02 & Hệ thống hiển thị các trạng thái loading, rỗng và lỗi của menu & Must & Mô phỏng API chậm, trả rỗng, trả lỗi thì UI hiện đúng từng trạng thái \\
FR-03 & Xem chi tiết sản phẩm; chọn size (S/M/L) và topping & Must & \texttt{GET /\allowbreak{}products/\allowbreak{}:id} trả đủ tùy chọn; giá cập nhật theo lựa chọn (S +0, M +5.000, L +10.000) \\
FR-04 & Thêm, sửa số lượng, xóa món; hiển thị tổng tiền & Must & Tổng tiền = $\Sigma$ (đơn giá + phụ phí size + topping) $\times$ số lượng \\
FR-05 & Đăng ký và đăng nhập; phát JWT; mật khẩu băm bằng bcrypt & Must & Đăng nhập đúng nhận JWT; sai mật khẩu trả 401; DB không lưu mật khẩu thô \\
FR-06 & Chặn đặt đơn khi chưa đăng nhập & Must & \texttt{POST /\allowbreak{}orders} không có JWT trả 401 \\
FR-07 & Tạo đơn từ giỏ hàng, trạng thái PENDING, trừ tồn kho trong cùng transaction & Must & Đơn tạo đúng; \texttt{stock} giảm; hết hàng thì từ chối và không tạo đơn \\
FR-08 & Thanh toán Ví hoặc Thẻ mô phỏng; thành công $\to$ PAID, thất bại $\to$ PAYMENT\_FAILED (API trả 200 kèm trạng thái đơn); thử lại được & Must & \texttt{POST /\allowbreak{}payments} đổi trạng thái đúng; thất bại giữ nguyên giỏ hàng \\
FR-09 & Gửi lại yêu cầu thanh toán với cùng Idempotency-Key không tạo thêm thanh toán/đơn & Must & Hai request cùng khóa cho một kết quả và một bản ghi \\
FR-10 & Máy trạng thái đơn; chặn chuyển sai trạng thái & Must & PENDING $\to$ PAID/PAYMENT\_FAILED/CANCELLED; PAID $\to$ PREPARING/CANCELLED; PREPARING $\to$ READY $\to$ COMPLETED; chuyển sai trả lỗi \\
FR-11 & Nhân viên cập nhật trạng thái đơn & Should & \texttt{PATCH /\allowbreak{}orders/\allowbreak{}:id/\allowbreak{}status} chỉ vai trò STAFF thực hiện được \\
FR-12 & Khách xem lịch sử và chi tiết đơn của mình & Must & \texttt{GET /\allowbreak{}orders/\allowbreak{}me} chỉ trả đơn của người dùng hiện tại; \texttt{GET /\allowbreak{}orders/\allowbreak{}:id} chỉ chủ đơn hoặc STAFF xem được (người khác nhận 404) \\
FR-13 & Đặt đồng thời không làm tồn kho âm; hủy đơn hoàn lại tồn kho & Must & Test chạy N request đồng thời, số đơn thành công $\leq$ tồn kho \\
FR-14 & Áp mã khuyến mãi hợp lệ để giảm giá & Should & Mã hợp lệ giảm đúng; mã sai/hết hạn bị từ chối \\
FR-15 & Cộng điểm khi đơn thanh toán thành công (PAID): 1 điểm / 10.000đ; hủy đơn đã PAID thì trừ lại; mỗi lần đổi điểm ghi một giao dịch điểm (\texttt{LoyaltyTransaction}) & Should & \texttt{User.points} tăng đúng số điểm khi đơn PAID và giảm lại khi hủy đơn đã PAID; có bản ghi \texttt{LoyaltyTransaction} loại EARN (và REVERSAL khi hủy) tương ứng \\
FR-16 & Giao diện quản lý sản phẩm và khuyến mãi & Could & Ngoài MVP \\
FR-17 & Kết nối cổng thanh toán thật & Won't & Không thực hiện \\
\end{longtable}}

\section{Yêu cầu phi chức năng}\label{sec:se-nfr}
{\small
\begin{longtable}{@{}p{1.3cm}p{5.3cm}p{1.2cm}p{6.6cm}@{}}
\caption{Yêu cầu phi chức năng và cách kiểm chứng}\label{tab:nfr}\\
\toprule \textbf{Mã} & \textbf{Yêu cầu} & \textbf{Ưu tiên} & \textbf{Tiêu chí chấp nhận} \\ \midrule \endfirsthead
\toprule \textbf{Mã} & \textbf{Yêu cầu} & \textbf{Ưu tiên} & \textbf{Tiêu chí chấp nhận} \\ \midrule \endhead
\bottomrule \endlastfoot
NFR-01 & \texttt{GET /\allowbreak{}products} phản hồi nhanh với dữ liệu seed & Should & Lần chậm thứ 19 trong 20 yêu cầu \texttt{GET /\allowbreak{}products} dưới 500 ms với dữ liệu mẫu (12 sản phẩm), đo cục bộ bằng Docker Compose. Kết quả ngày 24/10/2026: 19,7 ms. Ghi chú production (Render Singapore, gói Free, đo từ Việt Nam, 30 yêu cầu): byte đầu tiên về sau 0,43 đến 0,47 giây ở 24 yêu cầu và 0,74 đến 0,78 giây ở 6 yêu cầu, trong đó khoảng 0,25 giây là độ trễ mạng một vòng \\
NFR-02 & Mật khẩu băm; JWT có hạn; kiểm tra phân quyền theo vai trò & Must & Test xác thực và phân quyền; không có mật khẩu thô trong DB hoặc log \\
NFR-03 & Secret không nằm trong repository hoặc log CI & Must & Quét repo/log không thấy secret; \texttt{.env} nằm trong \texttt{.gitignore} \\
NFR-04 & Validate dữ liệu đầu vào ở cả frontend và backend & Must & Dữ liệu sai trả lỗi 4xx rõ ràng \\
NFR-05 & Giao diện mobile-first, dùng được trên màn hình điện thoại & Must & Kiểm tra trên viewport 360–414 px, không cuộn ngang \\
NFR-06 & Thông báo lỗi dễ hiểu, không lộ chi tiết kỹ thuật & Should & Lỗi hiển thị bằng tiếng Việt, không in stack trace \\
NFR-07 & Chạy toàn bộ hệ thống bằng \texttt{docker compose up} & Must & Người ngoài nhóm chạy được theo README \\
NFR-08 & CI chạy lint, test, build cho mọi PR; pipeline CD triển khai và smoke test & Must & Có log pipeline thành công và một lần phát hiện lỗi \\
NFR-09 & \texttt{GET /\allowbreak{}health} và log để vận hành & Must & Health trả \texttt{status: ok}; smoke test gọi health sau deploy \\
NFR-10 & Migration và seed tái tạo được database & Must & Chạy migration + seed trên DB trống cho kết quả giống nhau \\
NFR-11 & Mã nguồn có lint, định dạng và test tự động & Should & CI chạy lint và test \\
\end{longtable}}
Kết quả kiểm chứng thực tế của từng NFR nằm ở \cref{ch:truyvet} và \cref{ch:kiemthu}; NFR-01 đo thủ công, NFR-03 mới dò sơ bộ.

\section{Use case và đặc tả use case}\label{sec:se-usecase}
Sơ đồ use-case tổng quát ở \cref{fig:uc-overview}; phạm vi UC-01 ở \cref{fig:uc-checkout}. UC-01 là luồng sinh doanh thu, sai ở đây gây mất tiền hoặc sai tồn kho, nên là use case trọng tâm.
\hinh{usecase-overview.png}{.78\linewidth}{Use-case tổng quan của BrewLite}{fig:uc-overview}
\nguon{nhóm vẽ bằng PlantUML (\path{docs/diagrams/usecase-overview.puml})}
\hinh{usecase-checkout.png}{.85\linewidth}{Use-case của phạm vi UC-01}{fig:uc-checkout}
\nguon{nhóm vẽ bằng PlantUML (\path{docs/diagrams/usecase-checkout.puml})}

\begin{table}[htbp]\centering
\caption{Đặc tả UC-01: Đặt và thanh toán đơn hàng}\label{tab:uc01}
\begin{tabular}{@{}p{3.1cm}p{10.2cm}@{}}\toprule
Mã, tên & UC-01, Đặt và thanh toán đơn hàng \\
Actor chính / phụ & Customer / Payment Service (mô phỏng) \\
Trigger & Khách bấm ``Đặt đơn'' trong giỏ hàng \\
Yêu cầu liên quan & FR-04 đến FR-10, FR-13 đến FR-15, NFR-02, NFR-04 \\
Tiền điều kiện & Khách có JWT hợp lệ; giỏ có ít nhất một món \\
Hậu điều kiện & Thành công: đơn PAID, tồn kho đã trừ, có bản ghi thanh toán. Thanh toán thất bại: PAYMENT\_FAILED, giữ giỏ, cho thử lại. Hết hàng hoặc lỗi: không tạo đơn, tồn kho không đổi \\ \bottomrule
\end{tabular}
\end{table}
\textbf{Luồng chính.} (1) Khách mở giỏ và bấm ``Đặt đơn''. (2) Giao diện gửi \texttt{POST /orders} kèm JWT. (3) API kiểm tra JWT và validate. (4) API mở transaction: trừ tồn kho từng món, tạo Order (PENDING) và OrderItem, áp khuyến mãi nếu có, commit. (5) API trả mã đơn; giao diện hiển thị thanh toán. (6) Khách chọn Ví hoặc Thẻ. (7) Giao diện gửi \texttt{POST /payments} kèm \texttt{Idempotency-Key}. (8) Payment Service trả thành công. (9) API ghi Payment, chuyển đơn sang PAID. (10) Giao diện hiển thị xác nhận.

\begin{longtable}{@{}p{1cm}p{5.6cm}p{6.7cm}@{}}
\caption{Luồng thay thế (A) và ngoại lệ (E) của UC-01}\label{tab:uc01-alt}\\
\toprule \textbf{Mã} & \textbf{Điều kiện} & \textbf{Xử lý} \\ \midrule \endfirsthead
\toprule \textbf{Mã} & \textbf{Điều kiện} & \textbf{Xử lý} \\ \midrule \endhead
\bottomrule \endlastfoot
A1 & Mã khuyến mãi hợp lệ & Giảm giá, cập nhật tổng tiền trước khi đặt đơn \\
A2 & Gửi lại thanh toán cùng \texttt{Idempotency-Key} & Trả kết quả cũ, không tạo thanh toán thứ hai \\
A3 & STAFF hủy đơn ở trạng thái cho phép & CANCELLED, hoàn tồn kho, trừ lại điểm thưởng \\
E1 & Chưa đăng nhập hoặc JWT hết hạn & 401, chuyển sang đăng nhập, giữ giỏ \\
E2 & Dữ liệu không hợp lệ & 400 kèm thông báo rõ ràng \\
E3 & Món hết hàng hoặc không đủ số lượng & Rollback, không tạo đơn, báo món nào hết \\
E4 & Mã khuyến mãi sai hoặc hết hạn & Từ chối mã, cho đặt đơn không mã \\
E5 & Payment Service trả thất bại & PAYMENT\_FAILED, giữ giỏ, cho thử lại \\
E6 & Chuyển trạng thái không hợp lệ & 400, đơn giữ nguyên trạng thái \\
\end{longtable}
